\section{Introduction}
\label{sec:intro}

\IEEEPARstart{F}{ederated} learning lets several institutions train one predictive model without exchanging records~\cite{mcmahan2017,kairouz2021}. In the cross-silo setting that matters for healthcare, the participants are a handful of hospitals, each holding tens of thousands of encounters, each behind its own firewall, and each connected to the aggregation server by a link that the institution's network team controls~\cite{rieke2020,dayan2021}. In that setting the learning job is one tenant among many on an uplink with a fixed capacity, a scheduled transfer window and an audit trail. Whether that link is an institutional WAN, a 5G/6G edge slice serving a mobile clinic, or a satellite backhaul serving a remote site, the operator does not ask whether the model converges; it asks how many bytes the job will push per window and whether it will fit~\cite{lim2020,niknam2020,chen2021joint}. This article takes that question as the primary constraint and asks how a federated job should spend a stated communication quota.

Two obstacles compound. The first is statistical heterogeneity: sites serve different populations and use different coding conventions, so the IID assumption behind FedAvg fails and client updates diverge~\cite{zhao2018,li2022niid}. The second is communication: a parameter vector per client per round, every round, for tens of rounds. Prior work attacks the second obstacle either by shrinking each message (quantisation, sparsification, sketching~\cite{sattler2020,rothchild2020}) or by sending it less often (local steps, adaptive synchronisation periods~\cite{stich2019,wang2019adacomm,chen2018lag}). A third family keeps part of the model local for personalisation and sends the rest every round~\cite{arivazhagan2019,collins2021,liang2020,oh2022,pillutla2022}. Layer-wise schedules, which synchronise different parameter blocks at different frequencies~\cite{lee2023fedlama}, sit between the second and third family and rest on an intuition that early layers encode slowly-changing population structure and can therefore be averaged less often.

\textbf{We show that this intuition is wrong for tabular clinical data.} A fixed three-tier schedule that synchronises the input projection every five rounds and the hidden blocks every two loses \num{diabetesmlpFixedVsFedperDiff}~AUROC against FedPer on the Diabetes 130-US benchmark, over ten seeds, after Holm correction (Section~\ref{sec:results}). The input projection is exactly where heterogeneous coding conventions must be reconciled, so starving it is expensive. The failure is diagnostic, not fatal. A blockwise local-SGD bound (Section~\ref{sec:analysis}) makes the cost explicit: a block synchronised every $K_g$ rounds contributes a drift term proportional to $K_g^2\Gamma_g$, where $\Gamma_g$ is the block's cross-client gradient divergence. Sparsity is cheap only where $\Gamma_g$ is small, and $\Gamma_g$ is an empirical quantity that depends on the data, the architecture and the training stage. No schedule fixed before training can be right in general.

This motivates \tierfed{}, which turns the bound into a controller. Given a quota $\rho$ on the mean fraction of the shared payload per round, the server measures a dispersion proxy $\Gh_g$ of every tier it receives and solves the integer programme
\[
\min_{K_g\in\{1,\dots,K_{\max}\}}\ \sum_g K_g^2\,\Gh_g \quad\text{s.t.}\quad \sum_g w_g/K_g \le \rho ,
\]
exactly, by enumeration, every round. Rounds in which no tier is due transmit nothing. The classifier head is never transmitted. The only operator-facing quantity is $\rho$, which is a bandwidth quota rather than a learning-rate-like hyper-parameter.

\textbf{Contributions.}
\begin{itemize}
\item \emph{A negative result with a mechanism.} Fixed multi-rate aggregation underperforms simple personalisation on heterogeneous tabular clinical data across three architectures, and per-tier divergence measurements explain why (Sections~\ref{sec:results-fixed}, \ref{sec:results-arch}).
\item \emph{A bound that names the governing quantity.} A blockwise local-SGD bound whose drift term is $\sum_g K_g^2\tau_{\max}^2(\Gamma_g+\sigma_g^2)$, with the local-step count made explicit and the assumptions, including what the controller can and cannot observe, stated plainly (Section~\ref{sec:analysis}).
\item \emph{\tierfed{}, a budget-constrained scheduler.} An exact, sub-millisecond solver for the resulting integer programme, driven by online measurements, with a private head, composable with any compressor, and exposing one network-side knob (Section~\ref{sec:method}).
\item \emph{Network-level and security-level evaluation.} Communication time, energy and window compliance under five link profiles including 5G cell-edge, a 6G edge slice and LEO satellite; update-inversion on the slices each method actually transmits; membership inference on deployed models; and utility under client-level DP with per-client clipping and an exact R\'enyi accountant (Sections~\ref{sec:security}, \ref{sec:results-net}, \ref{sec:results-privacy}).
\item \emph{A reporting discipline.} Ten seeds, Holm-corrected paired tests, one-sided non-inferiority on both AUROC and AUPRC, bit-identical reruns, baselines on a single shared architecture table, and every number generated from released logs. Configurations in which \tierfed{} does not help are reported in the same tables.
\end{itemize}

The rest of the article is organised as follows. Section~\ref{sec:related} positions the work. Section~\ref{sec:system} states the network and learning model. Section~\ref{sec:method} presents \tierfed{}. Section~\ref{sec:analysis} gives the bound. Section~\ref{sec:security} states the threat model. Section~\ref{sec:setup} describes the data, protocol and statistics. Section~\ref{sec:results} reports results, Section~\ref{sec:discussion} discusses limitations, and Section~\ref{sec:conclusion} concludes.
